\reading{CR-11}{Portfolio Credit Risk}{STRIKE FIRST}{6}

\begin{crkeybox}[Learning objectives for this reading]
\begin{enumerate}[label=\textbf{\alph*.},leftmargin=2em]
\item Define and calculate default correlation for credit portfolios.
\item Identify drawbacks in using the correlation-based credit portfolio framework.
\item Assess the impact of correlation on a credit portfolio and its Credit VaR.
\item Describe the use of a single factor model to measure portfolio credit risk, including the impact of correlation.
\item Define beta and calculate the asset return correlation of any pair of firms using the single factor model.
\item Estimate the probability of a joint default of any pair of credits and the default correlation between any pair of credits using the single factor model.
\item Describe how Credit VaR can be calculated using a simulation of joint defaults.
\item Assess the effect of granularity on Credit VaR.
\end{enumerate}
\end{crkeybox}

\begin{crtrapbox}[Single-layer reading]
CR-11 exists only in the Crash layer. There is no Class Lecture treatment.
\end{crtrapbox}

% =====================================================================
\lo{CR-11 a}{Define and calculate default correlation for credit portfolios.}

The default correlation for a two-credit portfolio, assuming the outcomes are
Bernoulli-distributed random variables, rests on four mutually exclusive events.

\begin{center}
\footnotesize
\begin{tabular}{L{4.0cm} c c c L{3.6cm}}
\toprule
\thead{Event} & \thead{$x_1$} & \thead{$x_2$} & \thead{$x_1x_2$} &
\thead{Default probability} \\
\midrule
Firm 1 defaults & 1 & 0 & 0 & $\pi_1 - \pi_{12}$ \\
Firm 2 defaults & 0 & 1 & 0 & $\pi_2 - \pi_{12}$ \\
Both default & 1 & 1 & 1 & $\pi_{12}$ \\
No default & 0 & 0 & 0 & $1 - \pi_1 - \pi_2 + \pi_{12}$ \\
\bottomrule
\end{tabular}
\end{center}
\figcap{The four events and their probabilities. Note that ``firm 1 defaults''
means firm 1 alone, so the joint probability $\pi_{12}$ is subtracted out of it;
the four probabilities sum to 1 by construction.}

\begin{crfmlbox}[Default correlation for a two-credit portfolio]
\[ \rho_{12} \;=\; \frac{\pi_{12} - \pi_1\pi_2}
   {\sqrt{\pi_1(1-\pi_1)}\;\sqrt{\pi_2(1-\pi_2)}} \]
\vk{$\pi_1, \pi_2$ = the individual (marginal) default probabilities;
$\pi_{12}$ = the joint default probability. The denominator is the product of the
two Bernoulli standard deviations. If the two default events are independent then
$\pi_{12} = \pi_1\pi_2$ and the numerator, and hence the correlation, is zero.}
\end{crfmlbox}

% =====================================================================
\lo{CR-11 b}{Identify drawbacks in using the correlation-based credit portfolio framework.}

\begin{enumerate}
\item \term{Too many calculations.} The number of calculations required to
analyse all possible default scenarios increases dramatically as the number of
firms in the portfolio grows.
\item \term{Limited applicability.} The framework struggles to model certain
features like guarantees, revolving credit agreements, and convertible bonds.
\item \term{Data limitations.} Estimating default correlations is difficult due
to the rarity of defaults and the impact of specific data horizons and
industries.
\end{enumerate}

% =====================================================================
\lo{CR-11 c}{Assess the impact of correlation on a credit portfolio and its Credit VaR.}

The effects of default, default correlation, and loss given default are important
determinants in measuring credit portfolio risk.

\begin{crfmlbox}[Credit VaR]
\[ 95\%\ \mathrm{CVaR} \;=\; \text{95th percentile of the unrecovered credit loss}
   \;-\; \text{Expected Loss} \]
\vk{The percentile figure is the worst case loss (WCL) at the chosen confidence
level. Credit VaR is the \emph{excess} of the worst case over the expected loss,
which is why it can come out negative when the confidence level is not demanding
enough to produce any default at all.}
\end{crfmlbox}

\begin{crexbox}[Example 1: perfect correlation, $\rho = 1$]
Given a portfolio with notional value of \$1,000,000 and 20 credit positions.
Each credit has a PD of 2\% and an RR of 0. Each credit position is an obligation
from the \emph{same obligor}, so that the credit portfolio has a default
correlation equal to 1. What is the credit value at risk at the 99\% and 95\%
confidence levels for this portfolio?
\tcblower
With $\rho = 1$ the portfolio behaves as if there is only one credit: either
everything defaults or nothing does, with probability 2\%.

\[ \mathrm{EL} = \$1{,}000{,}000 \times 2\% = \$20{,}000 \]

\textbf{At 99\%.} The 1\% tail is inside the 2\% default probability, so the
worst case is total loss:
\[ \mathrm{WCL}(99\%) = \$1{,}000{,}000 \]
\[ \mathrm{CVaR}(99\%) = \$1{,}000{,}000 - \$20{,}000 = \$980{,}000 \]

\textbf{At 95\%.} The 5\% tail is wider than the 2\% default probability, so at
the 95th percentile no default has occurred:
\[ \mathrm{WCL}(95\%) = \$0 \]
\[ \mathrm{CVaR}(95\%) = \$0 - \$20{,}000 = -\$20{,}000 \]
\end{crexbox}

\begin{crexbox}[Example 2: zero correlation, $\rho = 0$]
Given a portfolio with a value of \$1,000,000 and 50 credits. Each credit is
equally weighted and has a terminal value of \$20,000 each if no default occurs.
Each credit has a PD of $\pi$ and an RR of zero. What is the credit VaR at the
95\% confidence level if $\pi$ is 2\% and the default correlation is 0? The 95th
percentile of the number of defaults based on this distribution is 3.
\tcblower
With $\rho = 0$ the number of defaults is binomially distributed.

\[ \mathrm{EL} = \$1{,}000{,}000 \times 2\% = \$20{,}000 \]
\[ \mathrm{WCL}(95\%) = 3 \times \$20{,}000 = \$60{,}000 \]
\[ \mathrm{CVaR}(95\%) = \$60{,}000 - \$20{,}000 = \$40{,}000 \]

\smallskip
Where the 3 comes from: under a binomial with $n=50$ and $p=0.02$, the cumulative
probability of 2 or fewer defaults is 92.16\% and of 3 or fewer is 98.22\%. Three
is therefore the smallest number of defaults that covers the 95th percentile.
\end{crexbox}

\begin{crkeybox}[What the pair of examples shows]
Same expected loss of \$20,000 in both. But at the 95\% level the perfectly
correlated portfolio has a credit VaR of $-\$20{,}000$ and the independent
portfolio has $+\$40{,}000$. Correlation does not move the expected loss; it moves
the shape of the tail, and it can move it in either direction depending on where
the confidence level falls relative to the default probability.
\end{crkeybox}

% =====================================================================
\lo{CR-11 d}{Describe the use of a single factor model to measure portfolio credit risk, including the impact of correlation.}

\subsection*{Model assumption}

The firm's asset return is represented as a function of two random variables: the
return on a \term{market factor} $m$ that captures the correlation between
default and the general state of the economy, and a \term{shock} $\varepsilon$
capturing idiosyncratic risk. Both $m$ and $\varepsilon$ are standard normal
variables and are not correlated with one another.

\begin{crfmlbox}[The single-factor decomposition]
\[ a_i \;=\; \beta_i m + \sqrt{1-\beta_i^{2}}\;\varepsilon_i
   \qquad\text{with}\qquad E(a_i) = 0, \quad \mathrm{Var}(a_i) = 1 \]
\vk{$a_i$ = the asset return of firm $i$; $\beta_i$ = the firm's loading on the
market factor. The coefficient on the shock is $\sqrt{1-\beta_i^{2}}$ precisely so
that the total variance stays at 1, which keeps $a_i$ standard normal.}
\end{crfmlbox}

\subsection*{Unconditional default probability}

The unconditional default distribution is a standard normal distribution. In the
single-factor model, the cumulative return distribution of any pair of credits
$i$ and $j$ is a \term{bivariate standard normal} with a correlation coefficient
equal to $\beta_i\beta_j$.

\subsection*{Conditional default probability}

Let $m$ take on a particular value $\bar{m}$. The distance to default now has only
one random driver, $\varepsilon$. The conditional distribution is normal with a
mean of $\beta_i\bar{m}$ and a standard deviation of $\sqrt{1-\beta_i^{2}}$.

\begin{crfmlbox}[Conditional cumulative default probability]
\[ p(\bar{m}) \;=\; \Phi\!\left(\frac{k_i - \beta_i\bar{m}}
   {\sqrt{1-\beta_i^{2}}}\right) \]
\vk{$k_i$ = the default threshold for firm $i$, the value of the asset return at
or below which the firm defaults. $\Phi$ is the standard normal cumulative
distribution function. The numerator is the \emph{new} distance to default given
the market factor; the denominator is the standard deviation under conditional
independence.}
\end{crfmlbox}

\begin{center}
\begin{tikzpicture}
\begin{axis}[
  width=12.4cm, height=6.0cm,
  xlabel={Market factor realization $\bar{m}$},
  ylabel={Conditional PD, $p(\bar{m})$},
  xlabel style={font=\sffamily\scriptsize}, ylabel style={font=\sffamily\scriptsize},
  tick label style={font=\sffamily\scriptsize},
  xmin=-4, xmax=2, ymin=0, ymax=1.0,
  grid=major, grid style={FRMRule!50, line width=0.3pt},
  legend style={font=\scriptsize\sffamily, draw=FRMRule, fill=white,
    at={(0.5,-0.30)}, anchor=north, legend columns=-1, column sep=10pt}]
\addplot[draw=FRMGreen, line width=1.1pt, smooth]
  coordinates {(-4.00,0.05920) (-3.50,0.04810) (-3.00,0.03873) (-2.50,0.03090)
  (-2.00,0.02443) (-1.50,0.01914) (-1.00,0.01486) (-0.50,0.01142) (0.00,0.00870)
  (0.50,0.00657) (1.00,0.00491) (1.50,0.00363) (2.00,0.00267)};
\addlegendentry{$\beta=0.2$}
\addplot[draw=FRMNavy, line width=1.1pt, smooth]
  coordinates {(-4.00,0.35158) (-3.50,0.25152) (-3.00,0.16893) (-2.50,0.10618)
  (-2.00,0.06230) (-1.50,0.03404) (-1.00,0.01730) (-0.50,0.00816) (0.00,0.00357)
  (0.50,0.00145) (1.00,0.00054) (1.50,0.00019) (2.00,0.00006)};
\addlegendentry{$\beta=0.5$}
\addplot[draw=FRMRed, line width=1.1pt, smooth]
  coordinates {(-4.00,0.92647) (-3.50,0.78328) (-3.00,0.54644) (-2.50,0.29116)
  (-2.00,0.11187) (-1.50,0.02983) (-1.00,0.00539) (-0.50,0.00065) (0.00,0.00005)
  (0.50,0.00000) (1.00,0.00000) (1.50,0.00000) (2.00,0.00000)};
\addlegendentry{$\beta=0.8$}
\addplot[only marks, mark=*, mark size=2pt, draw=FRMGold, fill=FRMGold]
  coordinates {(-0.5,0.00816)};
\node[font=\sffamily\scriptsize\bfseries, text=FRMGold, fill=white, inner sep=1.5pt]
  at (axis cs:0.78,0.14) {worked example: $0.82\%$};
\draw[-{Stealth[length=4pt]}, FRMGold] (axis cs:0.05,0.115) -- (axis cs:-0.45,0.03);
\end{axis}
\end{tikzpicture}
\end{center}
\figcap{Conditional default probability against the market factor, for a default
threshold of $k=-2.33$. All three curves describe firms with the same 1\%
unconditional default probability; what differs is $\beta$. A low-beta firm is
nearly flat, so the economy barely matters to it. A high-beta firm is close to
zero in good states and above 50\% by $\bar{m}=-3$. That fanning out \emph{is}
the correlation: it is what makes defaults arrive together.}

% =====================================================================
\lo{CR-11 e}{Define beta and calculate the asset return correlation of any pair of firms using the single factor model.}

\term{Beta} is the firm's loading on the market factor in the decomposition
above: the coefficient $\beta_i$ that determines how much of the firm's asset
return is driven by the systematic factor rather than by its own idiosyncratic
shock.

\begin{crfmlbox}[Asset return correlation of a pair]
\[ \mathrm{corr}(a_i, a_j) \;=\; \beta_i\beta_j \]
\vk{And where all firms share the same loading $\beta_i = \beta$, the pairwise
asset return correlation for any two firms is simply $\beta^{2}$.}
\end{crfmlbox}

\begin{crexbox}[Conditional default probability in the single-factor model]
Given a firm with a beta of 0.5 and a default threshold of $-2.33$, the
unconditional PD is $N(-2.33) = 1\%$. If the market factor is $-0.5$, what is the
conditional cumulative PD?
\tcblower
\textbf{Conditional mean.}
\[ \beta_i\bar{m} = 0.5 \times (-0.5) = -0.25 \]

\textbf{Conditional variance and standard deviation.}
\[ 1 - 0.5^{2} = 0.75 \qquad\Rightarrow\qquad \sqrt{0.75} = 0.8660 \]

\textbf{Conditional cumulative PD.}
\[ N\!\left[\frac{-2.33 - (-0.25)}{0.8660}\right]
   = N(-2.4018) = 0.8156\% \]

\smallskip
Note the direction: a market factor of $-0.5$ is a \emph{mildly good} state
relative to the default threshold, and the conditional PD of 0.82\% comes out
\emph{below} the unconditional 1\%.
\end{crexbox}

% =====================================================================
\lo{CR-11 f}{Estimate the probability of a joint default of any pair of credits and the default correlation between any pair of credits using the single factor model.}

\gapbadge

{\footnotesize\color{FRMGrey}The source notes jump from objective e to objective
g. This objective is absent from both layers. The following is written from the
GARP curriculum.}

\begin{crgapbox}[What the objective asks for]
Assume the parameters are the same for all firms, so $\beta_i = \beta$, $k_i = k$
and $\pi_i = \pi$. The pairwise asset return correlation for any two firms is
then $\beta^{2}$.

\smallskip
\textbf{Joint default probability.} Two firms default together when \emph{both}
asset returns fall at or below the common threshold $k$. Because the pair of
returns is bivariate standard normal with correlation $\beta^{2}$, the joint
default probability is the bivariate normal cumulative probability evaluated at
that corner:
\[ \pi_{12} \;=\; \Phi_2\bigl(k,\,k;\,\beta^{2}\bigr)
   \;=\; P\bigl[-\infty \leq a_1 \leq k,\; -\infty \leq a_2 \leq k\bigr] \]

\textbf{Default correlation.} Substituting this joint probability into the
two-credit correlation formula of objective a, with identical parameters, gives
\[ \rho \;=\; \frac{\Phi_2(k,k;\beta^{2}) - \pi^{2}}{\pi(1-\pi)} \]
which is the general expression from objective a specialised to a pair of
identical firms.
\end{crgapbox}

\begin{crkeybox}[Two sanity checks on the formula]
At $\beta = 0$ the asset returns are independent, so
$\Phi_2(k,k;0) = \pi^{2}$, the numerator vanishes, and the default correlation is
zero, as it must be.

\smallskip
\textbf{Asset correlation and default correlation are not the same number, and the
gap is large.} For a typical low investment-grade default probability of
$\pi = 0.01$, a default correlation of 0.05 requires an asset return correlation
$\beta^{2}$ of about 0.31, which is $\beta \approx 0.56$. Recovering $\beta$ from
a target default correlation requires a numerical procedure; there is no closed
form.
\end{crkeybox}

% =====================================================================
\lo{CR-11 g}{Describe how Credit VaR can be calculated using a simulation of joint defaults.}

A credit VaR under the copula methodology is computed by:

\begin{enumerate}
\item defining the copula function;
\item simulating default times;
\item obtaining market values and profit and loss data for each scenario using
the simulated default times; and
\item computing the portfolio distribution statistics by adding the simulated
terminal value results.
\end{enumerate}

% =====================================================================
\lo{CR-11 h}{Assess the effect of granularity on Credit VaR.}

\begin{crkeybox}[Granularity]
As a credit portfolio becomes \term{more granular}, the credit VaR
\term{decreases}.

\smallskip
However, when the default probability is \emph{low}, the credit VaR is not
impacted as much when the portfolio becomes more granular.
\end{crkeybox}
